% die Hauptdatei
% nur diese wird kompilliert
% der Rest wird hier hineinimportiert

\documentclass[a4paper, titlepage]{scrartcl}
\usepackage[utf8]{inputenc}
\usepackage[T1]{fontenc}
\usepackage[ngerman]{babel}
\usepackage{amsmath}
\usepackage{amsfonts}
\usepackage{graphicx}
\usepackage{ntheorem}
%\usepackage{amsthm}
%\usepackage{titlepic}
\usepackage{float}
\usepackage{enumerate}

\usepackage[scaled]{helvet}
\renewcommand*\familydefault{\sfdefault} %% Only if the base font of the document is to be sans serif

% verlinktes inhaltsverzeichnis (und mehr)
\usepackage[colorlinks=true,linkcolor=blue,pdfpagelabels=true]{hyperref}

\newcommand{\definition}[1]{\vspace{.5cm} \noindent \begin{large}\textbf{Def.: #1}\end{large}}
\newcommand{\beweis}[1]{\vspace{.5cm} \noindent \begin{large}\textbf{Bew.: #1}\end{large}}
\newcommand{\beispiel}[1]{\vspace{.5cm} \noindent \begin{large}\textbf{Bsp.: #1}\end{large}}
\newcommand{\bemerkung}[1]{\vspace{.5cm} \noindent \begin{large}\textbf{Bem.: #1}\end{large}}
\newcommand{\folgerung}[1]{\vspace{.5cm} \noindent \begin{large}\textbf{Folgerung #1:}\end{large}}
\newcommand{\behauptung}[1]{\vspace{.5cm} \noindent \begin{large}\textbf{Beh.: #1}\end{large}}
\newcommand{\lemma}[1]{\vspace{.5cm} \noindent \begin{large}\textbf{Lemma #1:}\end{large}}
\newcommand{\satz}[1]{\vspace{.5cm} \noindent \begin{large}\textbf{Satz #1:}\end{large}}

\newcommand{\widerspruch}{Widerspruch }

\newcommand{\datum}[1]{\begin{flushright}\textit{Vorlesung vom #1}\end{flushright}}

\newcommand{\AND}{\wedge}
\newcommand{\OR}{\vee}

\newenvironment{rcases}{%
  \left.\renewcommand*\lbrace.%
  \begin{cases}}%
{\end{cases}\right\rbrace}

\usepackage{yfonts}
\begin{document}
%\swabfamily

%\titlepic{\includegraphics[width=0.5\textwidth]{logo.png}}
%\title{Inoffizielles Vorlesungsskript zur Vorlesung „Logik“, WS 2013/2014}
%\author{gehalten von Prof. Volkmar Welker}
%\maketitle

\begin{figure}[h]
\centering
\includegraphics[width=0.6\textwidth]{vl_data/logo.png}
\end{figure}

\begin{center}
\begin{huge}inoffizielles Vorlesungsskript zur Vorlesung „Logik“, WS 2013/2014\end{huge}\\
\vspace{1cm}
\begin{large}gehalten von Prof. Volkmar Welker\end{large}\\
\vspace{1cm}
\begin{Large}\today \end{Large}
\end{center}

\noindent
\begin{center}
Die offizielle Vorlesungsseite:\\
\href{http://www.mathematik.uni-marburg.de/~welker/logik\_ws2013.html}{http://www.mathematik.uni-marburg.de/$\sim$welker/logik\_ws2013.html}
\end{center}

\noindent\\
\begin{center}\begin{Large}Fehler und Verbesserungsvorschläge bitte melden an: \href{mailto:beckers4@mathematik.uni-marburg.de}{beckers4@mathematik.uni-marburg.de}\end{Large}\end{center}
\vspace{1cm}
\begin{center}
empfohlenes Buch zur Vorlesung:\\
„Schöning, Logik für Informatiker“
\end{center}

\newpage

\tableofcontents

\newpage

% VL vom 14.10.2013

\section{Aussagenlogik}

\begin{flushright}
\textit{Vorlesung vom 14.10.2013}
\end{flushright}

Logische Systeme haben immer zwei Teile:
\begin{itemize}
\item Syntax (Formeln)
\item Semantik (Wahrheitswerte der Formeln)
\end{itemize}

\definition{(Syntax der Aussagenlogik)}
\begin{itemize}
\item Menge $A = \{ A_1, A_2, ... \}$ von atomaren Formeln
\item Wir definieren induktiv die Menge der Formeln der Aussagenlogik
    \begin{itemize}
        \item jede atomare Formel ist Formel
        \item sind $G$ und $H$ Formel, so auch \\
$( G \AND H )$ „$G$ und $H$“ \\
$(G \OR H)$ „$G$ oder $H$“ \\
$\neg G$ „nicht $G$“
    \end{itemize}
\end{itemize}

\bemerkung{}
Die Menge der Formeln der Aussagenlogik ist Sprache (= Menge von Wörtern) über Alphabet $A \cup \{\AND, \OR, \neg, (, ) \}$\\

\beispiel{}
\begin{itemize}
\item $ \left( ( A_{17} \OR A_2 ) \AND \neg (A_3 \AND A_4) \right)$\\
    Formel der Aussagenlogik
\item reine syntaktische Objekte\\
    $\neg \neg A \neq A$
\end{itemize}

\beweis{}
Sei $G$ eine Formel der Aussagenlogik\\
$T(G)$ die Menge der Teilformeln von $G$ induktiv definiert als\\
$T(G) = \{ G \}$ falls $G$ atomare Formel\\
$T(G) = T(G_1) \cup T(G_2) \cup \{G\}$, falls $G = (G_1 \OR G_2)$ oder $G=(G_1 \AND G_2)$\\
$T(G)=T(G_1) \cup \{G\}$\\

\beispiel{}
$G= \left( (A_{17} \OR A_2) \OR \neg (A_3 \AND A_4) \right)$\\
$T(G) = T \left( (A_{17} \OR A_2) \right) \cup T \left(\neg ( A_3 \AND A_4) \right) \cup \{G\}$\\
$= \{ A_{17}, A_2, (A_{17} \OR A_2) \} \cup T((A_3 \AND A_4)) \cup \{ \neg (A_3 \AND A_4) \} \cup \{G\}$\\
$= \{A_{17}, A_2, (A_{17} \OR A_2), A_3, A_4, (A_3 \AND A_4), \neg (A_3 \AND A_4), \left( (A_{17} \OR A_2) \AND \neg (A_3 \AND A_4) \right) \}$\\
    
\bemerkung{}
$T(G)$ ist die Menge der Formeln, die bei der induktiven Konstruktion von $G$ auftauchen.\\
\noindent\\
Sprechweisen:
%\begin{tabbing}
\begin{itemize}
\item Für $(G \AND H)$ sagt man auch „Konjunktion von $G$ und $H$“
\item Für $(G \OR H)$ sagt man auch „Disjunktion von $G$ und $H$“
\item Für $\neg G$ sagt man auch „Negation von G“.
\end{itemize}
%\end{tabbing}

\vspace{1cm}
\noindent
Abkürzende Schreibweisen:\\
\noindent\\
$G, H$ Formeln der Aussagenlogik
\begin{itemize}
\item $(G \rightarrow H)$ aus $G$ folgt $H$, für $(\neg G \AND H)$; $G$ impliziert $H$; Implikation
\item $(G \leftrightarrow H)$ $G$ äquivalent zu $H$, für $(G \rightarrow H) \AND (H \rightarrow G)$; Äquivalenz
\item $( \bigvee_{i=1}^{n} G_i )$ für $\left( … (G_1 \OR G_2) \OR G_3 ) … \OR G_n \right)$ \hspace{1cm}mit $G_1, …, G_n$ Formeln
\item $( \bigvee_{i=1}^n G_i )$ für $\left(… ( G_1 \AND G_2) \AND G_3) … \AND G_n \right)$ \hspace{1cm}mit $G_1, …, G_n$ Formeln
\end{itemize}

\vspace{1cm}
\definition{Semantik der Aussagenlogik}
\begin{itemize}
\item Sei $\emptyset \neq A' \subseteq A$ eine Teilmenge der atomaren Formeln
\item Eine Abbildung $f: A' \rightarrow \{W,F\}$ \hspace{1cm} (wahr, falsch)\\
    heißt Interpretation von $A'$
\item Eine Formel $G$ heißt Formel über $A'$ falls $T(G) \cap A \subseteq A'$
\item eine Interpretation $f: A \rightarrow \{W, F\}$ heißt passend zu Formel $G$, falls $G$ Formel über $A'$ ist.
\item Sei $f: A' \rightarrow \{W, F\}$ eine zur Formel $G$ passende Interpretation. Dann definieren wir $f(G)$ induktiv.\\
    $f(G) = f(( G_1 \AND G_2 )) = \begin{cases} W\ falls\ f(G_1) = f(G_2) = W \\ F\ sonst\end{cases}$\\
    $f(G) = f((G_1 \OR G_2)) = \begin{cases} W\ sonst \\ F\ falls\ f(G_1)=f(G_2)=F \end{cases}$\\
    $f(G) = f(\neg G_1) = \begin{cases} W\ falls\ f(G_1)=F \\ F\ falls\ f(G_1)=W \end{cases}$\\
    für $G=\neg G_1$
\end{itemize}

\beispiel{Semantik der abkürzenden Schreibweisen}

\begin{tabular}{c|c|c|c}
$G$ & $H$ & $(G \rightarrow H)$ & $(G \leftrightarrow H)$\\
\hline
F & F & W & W\\
F & W & W & F\\
W & F & F & F\\
W & W & W & W\\
\end{tabular}

\lemma{1.1}
Sei $A'$ eine Menge von $n$ atomaren Formeln. Dann gibt es $2^n$ Interpretationen von $A'$\\

\beweis{}
Für jede atomare Formel mit $A'$ gibt es $2$ Möglichkeiten für das Bild unter $f$\\
$\rightarrow$ Die atomaren Formeln $A'$ können unabhängig voneinander interpretiert werden\\
$\Rightarrow$ $\#$ Interpretationen $=\underbrace{2 \cdot 2 \cdot 2}_n = 2^n$ \hspace{1cm}\ensuremath{\Box}\\

\definition{}
Sei $G$ eine Formel über $A' = \{A_1, …, A_n \}$\\
seien $f_1, …, f_{2^n}$ die $2^n$ Interpretationen von $A'$\\
Dann heißt das Schema\\
\begin{tabular}{c|c|c|c}
$A_1$ & … & $A_n$ & $G$\\
\hline
$f_1 (A_1)$ & … & $f_1 (A_n)$ & $f_1 (G)$\\
\vdots & & \vdots & \\
$f_{2^n} (A_1)$ & … & $f_{2^n} (A_n)$ & $f_{2^n} (G)$
\end{tabular}\\
Wahrheitstabelle von $G$\\

\beispiel{}
$(A_1 \OR (A_2 \AND \neg A_3))$\\

\begin{tabular}{c|c|c|c}
$A_1$ & $A_2$ & $A_3$ & $(A_1 \OR (A_2 \AND \neg A_3))$\\
\hline
F & F & F & F\\
F & F & W & F\\
F & W & F & W\\
W & F & F & W\\
F & W & W & F\\
W & F & W & W\\
W & W & F & W\\
W & W & W & W\\
\end{tabular}\\
    
\definition{}
Sei $F_{A'}$ die Menge aller Formeln über $A'$\\
Für $G, H \in F_{A'}$ sagen wir $G$ ist semantisch äquivalent zu $H$, falls $f(G) = f(H)$ für alle Interpretationen $f:A' \rightarrow \{W,F\}$\\
Wir schreiben $G \equiv H$\\

\datum{17.10.2013}

\definition{}
Sei $F_n$ die Menge aller Formeln über $\{A_1, …, A_n\}$ eine Teilmenge $K \subseteq F_n$ heißt semantische Klasse, falls $G \equiv H$ für $G, H \in K$ und $G \not\equiv H$ für $G\in K$ \hspace{.5cm} $H \in F_n \backslash K$\\

\beispiel{}
$F_1 = \{A_1, \neg A_1, A_1 \OR A_1, A_1 \AND A_1, … \}$\\

\begin{minipage}{0.15\textwidth}
\begin{tabular}{c|c}
$A_1$ & $A_1$\\
\hline
F & F\\
W & W\\
\end{tabular}
\end{minipage}
\hfill
\begin{minipage}{0.15\textwidth}
\begin{tabular}{c|c}
$A_1$ & $\neg A_1$ \\
\hline
F & W\\
W & F\\
\end{tabular}
\end{minipage}
\hfill
\begin{minipage}{0.15\textwidth}
\begin{tabular}{c|c}
$A_1$ & $A_1 \OR A_1$\\
\hline
F & F\\
W & W\\
\end{tabular}
\end{minipage}
\hfill
\begin{minipage}{0.15\textwidth}
\begin{tabular}{c|c}
$A_1$ & $A_1 \OR \neg A_1$\\
\hline
F & W\\
W & F\\
\end{tabular}
\end{minipage}
\hfill
\begin{minipage}{0.15\textwidth}
\begin{tabular}{c|c}
$A_1$ & $A$\\
\hline
F & F\\
W & F\\
\end{tabular}
\end{minipage}\\
\noindent\\
$A_1 \equiv A_1 \OR A_1$\\
\noindent\\
Das sind alle Wahrheitstabellen, die bei einer atomaren Formel möglich sind.\\
\noindent\\
Die semantischen Klassen sind:\\
$K_1 = \{G \in F_1 | G \equiv A_1\}$\\
$K_2 = \{G \in F_1 | G \equiv \neg A_1\}$\\
$K_3 = \{G \in F_1 | G \equiv A_1 \OR \neg A_1\}$\\
$K_4 = \{G \in F_1 | G \equiv A_1 \AND \neg A_1\}$\\
    
\bemerkung{}
„$\equiv$“ ist Äquivalenzrelation auf $F_n$ und die semantischen Klassen sind die Äquivalenzklasse bzgl „$\equiv$“\\
    
\bemerkung{}
\begin{itemize}
\item Die Elemente einer semantischen Klasse sind alle Formeln mit der gleichen Wahrheitstabellen
\item Jede semantische Klasse enthält unendlich viele Formeln\\
    $(G \equiv \neg \neg G \equiv \neg \neg \neg \neg G \equiv … )$
\end{itemize}

\definition{}
\begin{itemize}
\item Eine Formel $G$ heißt gültig oder Tautologie, falls $f(G) = W$ für jede Interpretation $f$
\item Eine Formel $G$ heißt erfüllbar, falls es eine Interpretation $f$ gibt mit $f(G) = W$
\item Eine Formel $G$ heißt nicht erfüllbar, unerfüllbar oder Antilogie, falls $f(G) = F$ für alle Interpretationen $f$
\end{itemize}

\beispiel{}\\
$\begin{rcases} A_1 \OR \neg A_1\ g\ddot{u}ltig \\ A_1 \\ \neg A_1 \end{rcases} erf\ddot{u}llbar$\\
$A_1 \AND \neg A_1$ unerfüllbar\\

\noindent\\
Ein zentrales Ziel der Vorlesung wird es sein, Algorithmen zu entwickeln, die zur Gegebenen Formel $G$ entscheiden, da diese erfüllbar ist.\\
    
\noindent\\
\bemerkung{}
\begin{enumerate}[i)]
\item $G$ gültig $\Rightarrow G$ ist er erfüllbar
\item $G$ ist unerfüllbar $\Leftrightarrow \neg G$ gültig
\item $G, H$ gültig $\Rightarrow G \equiv H$
\item $G, H$ nicht erfüllbar $\Rightarrow G \equiv H$
\end{enumerate}

\noindent
Wir untersuchen nun die Ausdruckskraft der Aussagenlogik:\\
Gegeben eine Wahrheitstabelle\\
z. B.\\
\begin{tabular}{c|c|c|c}
$A_1$ & $A_2$ & $A_3$ & ?\\
\hline
F & F & F & W\\
F & F & W & F\\
F & W & F & F\\
W & F & F & W\\
F & W & W & F\\
W & F & W & F\\
W & W & F & W\\
W & W & W & F
\end{tabular}\\
Gibt es Formel $G$ mit dieser Wahrheitstabelle?\\

\lemma{1.2}
Über eine Menge von $n$ atomaren Formeln gibt es $2^{2^{n}}$ semantische Klassen. Insbesondere gibt es zu jeder Wahrheitstabelle eine Formel mit dieser Wahrheitstabelle.\\
    
\beweis{}
$\rightarrow$ Zeige: Es gibt $2^{2^n}$ Wahrheitstabellen über $n$ atomare Formeln. Die $2^n$ Interpretationen (Lemma 1.1) legen die Wahrheitstabelle bis auf die letzte Spalte fest.\\
Für jede Zeile gibt es zwei Möglichkeiten für den Eintrag in die letzte Spalte\\
\noindent\\
$\underbrace{2 \cdot 2 \cdot … \cdot 2}_{2^n} = 2^{2^n}$ Möglichkeiten\\
    
\noindent\\
$\rightarrow$ Zeige: Zu jeder Wahrheitstabelle gibt es Formel mit dieser Wahrheitstabelle\\
    
\noindent\\
Konstruktion\\

\begin{tabular}{c|c|c|c}
$A_1$ &  & $A_n$ &  \\
\hline
$f_1 F$ & … & $F$ & $W_1$\\
\vdots   &  & \vdots & \vdots \\
$f_n W$ & … & $W$ & $W_{2^n}$ \\
\end{tabular}
\hspace{1cm}$W_1, …, W_{2^n} \in \{W, F\}$\\
    
\noindent\\
Seien $i_1, …, i_l$ die Indizes mit $W_{i_j} = W$ für $1 \leq j \leq l$\\
$G_k = \bigwedge_{m=1}^n H_{k,m}$, $H_{k,m} = \begin{cases} A_m\ falls\ f_{i_k}\ (A_m)=W \\ \neg A_m\ falls\ f_{i_k}\ (A_m)=F \end{cases}$\\
    
\noindent\\
Setz $G = \bigvee_{k=1}^l G_k$\\
    
\noindent\\
Sei $f$ Interpretation\\
\begin{tabbing}
$G(G_k) = W$ \= $\Leftrightarrow f(H_{2,m}) = W\ f\ddot{u}r\ m=1, …, m$\\
\> $\Leftrightarrow f(f_m) = W\ f\ddot{u}r\ W\ mit\ f_{i,k}(A_m) = W$\\
\> = $f(\neg A_m) = W\ f\ddot{u}r\ W\ mit\ f_{i,k}(A_m) = F$\\
\> $\Leftrightarrow f = f_{i,k}$\\
\end{tabbing}

\noindent\\
\begin{tabbing}
$f(G) = W$ \= $\Leftrightarrow \exists k : f(G_k) = W$\\
\> $\Leftrightarrow \exists k: f= f_{i,k}$\\
\end{tabbing}
$\Rightarrow G$ hat die gegebene Wahrheitstabelle\\

\beispiel{}\\
\begin{tabular}{cc|c|c|c}
 &  & $A_1$ & $A_2$ &   \\
\hline
$\rightarrow$ & $f_1$ & F & F & \textcircled{W}\\
$\rightarrow$ & $f_2$ & W & F & \textcircled{W}\\
 & $f_3$ & F & W & F\\
 & $f_4$ & W & W & F\\
\end{tabular}\\
(suche Zeilen, die hinten $W$ sind)\\

\noindent
$i_1 = 1, i_2 = 2, l=2$\\
$G_1 = (\neg A_1 \AND \neg A_2)$\\
$G_2 = (A_1 \AND \neg A_2)$\\
$G = G_1 \OR G_2 = (( \neg A_1 \AND \neg A_2) \OR (A_1 \AND \neg A))$\\

\vspace{1cm}
\noindent
Regeln zur semantischen Äquivalenz.

\satz{1.3} Seien $G, H, I$ Formeln der Aussagenlogik\\
Dann gilt
\begin{tabbing}
Idempotenz \hspace{1cm} \= $(G \AND G) \equiv G$\\
\> $(G \OR H) \equiv (H \OR G)$\\
%\end{tabbing}
\\
%\begin{tabbing}
Assoziativität \> $((G \AND H) \AND I ) \equiv (G \AND (H \AND I))$\\
\> $((G \OR H) \OR I) \equiv (G \OR (H \OR I ))$\\
%\end{tabbing}
\\
%\begin{tabbing}
Distributivität \> $(G \AND ( H \OR I )) \equiv ((G \AND H ) \OR (G \AND I ))$\\
\> $(G \OR ( H \AND I )) \equiv (( G \OR H ) \AND ( G \OR I ))$\\
%\end{tabbing}
\\
%\begin{tabbing}
Absorption \> $(G \AND ( G \OR H )) \equiv G$\\
\> $(G \OR ( G \AND H )) \equiv G$\\
%\end{tabbing}
\\
Doppelnegation \> $\neg \neg G \equiv G$\\
\\
%\begin{tabbing}
deMorgan Regel \> $\neg ( G \AND H ) \equiv \neg G \OR \neg H$\\
\> $\neg ( G \OR H) \equiv \neg G \AND \neg H$
\end{tabbing}

\beweis{}
Beispielhaft Absorption

\begin{tabular}{c|c|c|c}
$G$ & $H$ & $G \AND ( G \OR H)$ & $G$ \\
\hline
F & F & F & F \\
F & W & F & F \\
W & F & W & W \\
W & W & W & W 
\end{tabular}
$\Rightarrow (G \AND ( G \OR H )) \equiv G$

\lemma{1.4}
$G, H, G', H'$ Formeln der Aussagenlogik\\
und $G \equiv G'$ und $H \equiv H'$\\
\noindent\\
$\Rightarrow$\\
$(G \AND H) \equiv (G' \AND H')$\\
$(G \OR H) \equiv (G' \OR H')$\\
$\neg G \equiv \neg G'$

\beweis{}\\
\begin{tabular}{c|c|c|c|c|c}
$G$ & $H$ & $G'$ & $H'$ & $G \AND H$ & $G' \AND H'$ \\
\hline
F & F & F & F & F & F \\
F & W & F & W & F & F \\
W & F & W & F & F & F \\
W & W & W & W & W & W \\
\end{tabular}
$\Rightarrow (G \AND H ) \equiv ( G' \AND H')$\\
\noindent\\
andere analog

\satz{1.5}
Seien $G, H, I$ Formeln der Aussagenlogik und $G \equiv H$ und $G$ ist Teilformel von $I$.\\
Sei $I'$ eine Formel, die aus $I$ entsteht indem man ein Vorkommen von $G$ durch $H$ ersetzt\\
Dann $I \equiv I'$
\datum{21.10.2013} 


\satz{1.5} $G \equiv H$\\
$I$ mit $G$ Teilformel\\
$I'$ ersetze ein Vorkommen von $G$ in $I$ durch $H$\\
$\Rightarrow I \equiv I'$

\beispiel{}\\
$\begin{rcases}G = \neg ( \neg A_1 \OR A_3)\\
H = (A_1 \AND \neg A_3) \end{rcases} G \equiv H$\\
\noindent\\
$I = ((( \neg A_{2} \OR A_4) \AND ( A_6 \OR \neg ( \neg A_1 \OR A_3))) \AND \neg ( \neg A_1 \OR A_3))$\\
$I' = ((( \neg A_2 \OR A_4) \AND (A_6 \OR ( A_1 \AND \neg A_3))) \AND \neg ( \neg A_1 \OR A_2))$\\

\beweis{}
Induktion über die Länge  $l = \#$ zeichen von $I$\\
I.A.: $l$ = Länge von $G$\\
$\Rightarrow I = G$ und $I' = H$\\
$\Rightarrow I = G \equiv H = I'$\\
\noindent\\
I.V.: Die Aussage gilt für alle Formeln der Länge $\leq l$\\
I.S.: Sei $I$ Formel der Länge $l+1$, die $G$ als Teilformel enthält\\
\noindent\\
\underline{1. Fall}\\
$I = (I_1 \AND I_2)$ mit Formeln $I_1$ und $I_2$ der Länge $\leq l$

\begin{enumerate}
\item[i)] $I' = (I'_1 \AND I_2)$\\
$I'_1$ entsteht aus $I_1$ durch ersetzen der Teilformel $G$ in $I$, durch $H$\\
Sei Länge von $H_1 \leq l$ folgt nach I.V\\
$I_1 = = I'_1$\\
Also\\
$\begin{rcases} I = (I_1 \AND I_2)$, $I'=(I'_1 \AND I_2) \\
I_1 \equiv I'_1 \\
I_2 = I_2 \end{rcases} I' \equiv I$\\

\item[ii)] $I' = I_1 \AND I'_2$\\
$I'_2$ entsteht aus $I_2$ durch ersetzen der Teilformel $G$ in $I_2$ durch $H$\\
weiter analog zu (i)
\end{enumerate}

\noindent
\underline{2. Fall}\\
$I = (I_1 \OR I_2)$\\
analog zu 1. Fall\\

\noindent
\underline{3. Fall}\\
$I_1 = \neg I_1$ ist Formel $I_1$ der länge $\leq l$\\
$\rightarrow I' = \neg I'_1$ und $I'_1$ entsteht aus $I_1$ durch ersetzen der Teilformel $G$ durch $H$\\

\noindent
Nach I.V. \hspace{1cm} $I_1 \equiv I'_1$\\
$\Rightarrow^{Lemma\ 1.4}$ (iii) $I = \neg I_1 \equiv \neg I'_1 = I'$ $\Box$\\
$\rightarrow$ Der Beweis von Satz 1.5 ist ein erstes Beispiel des Beweisprinzips „Induktion über Formelaufbau“\\
(„Strukturelle Induktion“)

\behauptung{„Aussage über Formeln“}

\beweis{}
Induktion über Formelaufbau\\
I.A.: Aussage für Kürzest mögliche Formel\\
I.S.: Die Aussage gilt für die 3 Möglichkeiten\\
für $(G \AND H)$, $(G \OR H)$, $\neg G$\\
falls die Aussage für $G$ und $H$ gilt.

\definition{}
Sei $\Sigma$ eine Menge von Formeln\\
$\Sigma$ heißt erfüllbar, falls es eine Interpretation $f$ gibt mit $f(G) = W\ \forall\ G \in \Sigma$\\
so ein $f$ heißt Modell $\Sigma$

\noindent
$\Sigma$ heißt unerfüllbar oder nicht erfüllbar, falls es eine Interpretation $f$ gibt mit $f(G) = W$

\noindent
für alle $G \in \Sigma$\\

\beispiel{}
$\Sigma = \{A_1, A_1 \AND A_2, A_1 \AND A_2 \AND A_3, … \}$\\
$= \{A_1 \AND … \AND A \ |\ n=1 \}$\\

\noindent
erfüllbar $f(A_i) = W$ für $i=1$ ist einziges Modell\\
$\Sigma = \{A_1 \AND A_2, \neg ( \neg A_1 \OR A_2) \}$ nicht erfüllbar\\

\beweis{}
\begin{enumerate}
\item[i)] $G$ erfüllbar $\Leftrightarrow \{G\}$ erfüllbar
\item[ii)] $\Sigma$ erfüllbar $\Rightarrow G$ erfüllbar für alle $G \in \Sigma$
\item[iii)] $G$ erfüllbar für alle $G \in \Sigma \not \Rightarrow \Sigma$ erfüllbar
\end{enumerate}
$\Sigma = \{A_1, \neg A_1 \}$ unerfüllbar\\
\noindent\\
$A_1, \neg A_1$ erfüllbar wenn $A_1 = W$ bzw. $F$

\lemma{1.6}
Sei $\Sigma = \{G_1, G_n \}$ endlich, dann folgt\\
$\Sigma$ erfüllbar $\leftrightarrow (G_1 \AND … \AND G_n)$ erfüllbar

\beweis{}
$\Sigma$ erfüllbar $\Leftrightarrow$ Es gibt Interpretation $f$ mit $f(G_1) = W$, $i=1, …, n$\\
$\Leftrightarrow$ Es gibt Interpretation $f$ mit $f(G_1 \AND … \AND G_n) = W$\\
$\Leftrightarrow ( G_1 \AND … \AND G_n)$ erfüllbar $\Box$

\definition{}
Eine Formelmenge $\Sigma$ impliziert semantisch eine Formel $H$, falls für jedes Modell $f$ von $\Sigma$ gilt $f(H) = W$; d.h. $f$ ist auch Modell von $\{H\}$. Wir schreiben $\Sigma \models H$. Ist $\Sigma = \{G\}$ so schreiben wir auch $G \models H $ für $G$ impliziert semantisch $H$.

\beispiel{}
$A_1 \AND A_2 \models A_1 \OR A_2$\\
$f(A_1 \AND A_2) = W \Rightarrow f(A_1) = f(A_2) = W$\\
$\Rightarrow f(A_1 \OR A_2) = W$\\
$A_1 \OR A_2 \not \models A_1 \AND A_2$\\

\noindent
für $f(A_1) = W,\ f(A_2) = F$ gilt $f(A_1 \OR A_2) = W$\\
aber $f(A_1 \AND A_2) = F$\\
falls $G \models H$ und $f(G) = F$ muss nicht $f(H) = f$ gelten.\\
$(A_1 \AND A_2) \models (A_1 \OR A_2)$\\
für $f(A_1) = W,\ f(A_2)=F$ gilt $f(A_1 \AND A_2) = F$\\
aber $f(A_1 \OR A_2 ) = W$

\bemerkung{}
$G, H$ Formeln\\
$G \equiv H \Leftrightarrow ( G \models H)$ und $(H \models G)$\\

\beweis{}
$G \equiv H$, für jede Interpretation $f$ gilt $f(H) = f(G)$\\
$\Leftrightarrow$ für jede Interpretation $f$ gilt\\
$f(G) = W \Leftrightarrow f(H) = W$\\
$\Leftrightarrow G \models H$ und $H \models G$ $\Box$

\beispiel{}
Es gibt Formel $G, H$ mit $G \not \models H, H \not \models G$\\
$G = A_1$, $H=\neg A_1$

\satz{1.7} \textbf{Gaig'scher Interpolationssatz}\\
Seien $G$ und $H$ Formel mit $G \models H$\\
Dann gilt mindestens eine der folgenden 3 Aussagen:
\begin{itemize}
\item[i)] $H$ ist gültig
\item[ii)] $G$ ist unerfüllbar
\item[iii)] Es gibt eine Formel $I$ mit $G \models I$, $I \models H$
\end{itemize}
und jede atomare Teilformel von $I$ ist atomare Teilformel von $G$ und $H$.

\beispiel{}
$G = (A_1 \AND A_2)$, $H = (A_2 \OR A_3)$\\
$G \models I$ \hspace{1cm} $I \models H$ \hspace{1cm} für $I= A_2$\\
$(A_1 \AND A_2) \models A_2$ \hspace{1cm} $A_2 \models (A-2 \OR A_3)$

\beweis{}
Es genügt zu Zeigen\\
$G \models H$ und $H$ nicht gültig mit $G$ erfüllbar $\Rightarrow$ (iii)\\

\noindent
Induktion über die Anzahl der atomaren Teilformeln von $G$, die nicht in $H$ vorkommen.\\
I.A: $n=0$ \hspace{1cm} Können $I = G$ wählen\\
$G \models I = G$ \hspace{1cm} $G = I \models H$\\


\datum{24.10.2013}


\satz{1.7} \textbf(Graig'scher Interpolationssatz)\\
Seien $G,H$ Formeln mit $G \models H$\\
Dann gilt mindestens eine der folgeden Aussagen:\\
\begin{enumerate}
\item[i)] $H$ gültig
\item[ii)] $G$ unerfüllbar
\item[iii)] Es gibt eie Formel $I$ mit den Eigenschaften $G \models I$, $I \models H$\\
und jede atomare Teilformel von $I$ ist Teilformel von $G$ und $H$.
\end{enumerate}

\beweis{}
Es genügt zu zeigen:\\
%$\begin{rcases}$
\begin{itemize}
\item $H$ nicht gültig
\item $G$ erfüllbar
\item $G \models H$
\end{itemize}
%$\end{rcases}$
$\Rightarrow$ iii)\\

\noindent
Induktion über die Anzahl der atomaren Teilformeln von $G$, die nicht in $H$ enthalten sind.\\

\noindent
I.A.: Sei $n=0$. Dann ist jede atomare Teilformel von $G$ in $H$ enthalten.\\
Wähle $I=G$, dann enthält $I$ nur atomare Teilformeln, die in $G$ und $H$ enthalten sind und $G \models I$ und $I \models H$.\\

\noindent
I.V.:
\begin{itemize}
\item $G \models H$ 
\item $H$ nicht gültig
\item $G$ erfüllbar
\item Die Anzahl der atomaren Teilformeln von $G$, die nicht in $H$ enthalten sind, ist $n$.
\end{itemize}
$\Rightarrow$ iii)\\

\noindent
I.S.: $(n \rightarrow n+1)$\\
\begin{itemize}
\item $G \models H$
\item $H$ nicht gültig
\item $G$ erfüllbar
\item Die Anzahl der atomaren Formeln von $G$, die nicht in $H$ enthalten sind, ist $n+1$
\end{itemize}
\noindent
Sei $A_i$ eine atomare Formel, die in $G$ aber nicht in $H$ enthalten ist.\\
Sei $G_W$/$G_F$ die Formel, die aus $G$ entsteht, wenn jedes Vorkommen von $A_i$ ersetzt wird durch eine gültige/ungültige Formel, die nur atomare Formeln von $G$ und $H$ verwendet.\\

\noindent
Setze $G' := (G_W \OR G_F)$
\begin{enumerate}
\item[1)] Es gibt eine atomare Formel, die in $G$ und in $H$ enthalten ist.\\
(d.h. $G'$ ist wohldefiniert)\\
\underline{Beweis:} Angenommen $G$ und $H$ besitzen keine gemeinsame atomare Formel.\\
Da $G$ erfüllbar ist, gibt es eine Interpretation $f_G$ der atomaren Teilformeln von $G$ mit $f_G (G) = W$\\
Da $H$ nicht gültig ist, gibt es eine Interpretation $f_H$ der atomaren Teilformeln von $H$ mit $f_H (H) = F$\\
Da die Menge der atomaren Teilformeln von $G$ disjunkt zur Menge der atomaren Teilformeln von $H$ ist, gibt es eine gemeinsame Erweiterung $f$ von $f_G$ und $f_H$\\
Damit ist $f$ eine Interpretation der atomaren Teilformeln von $G$ und $H$ mit:
\begin{itemize}
\item $f(G) = f_G(G) = W$
\item $f(H) = f_H(H) = F$\\ \widerspruch zu $G \models H$
\end{itemize}
$\Rightarrow$ Ann. falsch $\Rightarrow$ Beh.

\item[2)] \underline{Beh:} $G' \models H$ \\
\underline{Beweis:} Es genügt zu zeigen:
\begin{enumerate}
\item[a)] $G_W \models H$
\item[b)] $G_F \models H$
\end{enumerate}

zu a)\\
Sei $f$ eine Interpretation der atomaren Teilformeln $G_W$ und $H$ mit $f(G_W) = W$.\\
Erweitere $f$ zu einer Interpretation $f'$ der atomaren Formeln von $G$ und $H$ durch\\
$f'(A_j) = \begin{cases} f(A_j),\ falls\ A_j \neq A_i \\ W,\ falls\ A_j = A_i \end{cases}$\\
$\Rightarrow f'(G) = f(G_W) = W$\\
$\Rightarrow^{(G \models H)} f(H) = f'(H) = W$\\
$\Rightarrow G_W \models H$\\

zu b)\\
Sei $f$ eine Interpretation von $G_F$ und $H$ mit $f(G_F) = W$.\\
Erweitere $f$ zu einer Interpretation $f'$ von $G$ und $H$ durch\\
$f'(A_j) = \begin{cases} f(A_j),\ falls\ A_j \neq A_i \\ F,\ falls\ A_j = A_i \end{cases}$\\
$\Rightarrow^{(G \models H)} f(H) = f'(H) = W$\\
$\Rightarrow G_F \models H$

\item[3)] \underline{Beh:} Die Anzahl der atomaren Formeln in $G'$, die nicht in $H$ enthalten sind ist $n$.\\
\underline{Beweis:} klar

\item[4)] \underline{Beh:} $G'$ ist erfüllbar.\\
\underline{Beweis:}\\
Sei $f$ eine Interpretation der atomaren Formeln in $G$ mit $f(G) = W$.\\
($G$ erfüllbar)\\
Entweder $f(A_i) = W$, d.h. $f(G_W) = W$\\
oder $f(A_i) = F$, d.h. $f(G_F) = W$\\
Somit ist $f(G') = f(G_W \OR G_F) = W$

\item[5)] \underline{Beh:} Es gibt ein $I$, wie in der Beh. des Satzes.\\
\underline{Beweis:} Nach I.V. gibt es eine Formel $I$, welche nur die atomare Formeln von $G'$ und $H$ verwendet mit $G' \models I$ und $I \models H$.

\begin{itemize}
\item $I$ verwendet nur die atomaren Formeln von $G$ und $H$
\item $I \models H$ 
\item zz: $G \models I$\\
Sei $f$ eine Interpretation der atomaren Aussagen von $G$ und $H$ mit $f(G) = W$\\
Dann gilt entweder\\
$f(A_i) = W$ und $f(G_W) = W$\\
oder\\
$f(A_i) = F$ und $f(G_F) = W$\\
Somit folgt $f(G_W \OR G_F) = W$\\
$\Rightarrow^{(G' \models H)} f(I) = W$\\
$\Rightarrow G \models H$ $\Box$
\end{itemize}

\end{enumerate}


\section{Normalformen}

\beispiel{}
$\rightarrow (A_1 \OR \neg A_2 ) \equiv ( \neg A_1 \AND \neg \neg A_2) \equiv ( \neg A_1 \AND A_2)$\\

\definition{}
Ein \underline{Literal} ist eine atomare Formel oder deren Negation. \\

\noindent
Sprechweisen:\\
$A_i \rightarrow$ atomare Formel $\rightarrow$ positives Literal\\
$\neg A_i \rightarrow$ Negation einer atomaren Formel $\rightarrow$ negatives Literal\\

\definition{} Sei $\{A_1, …, A_n\}$ eine Menge von atomaren Formeln\\

\noindent
Ein \underline{Minterm} über $\{A_1, …, A_n\}$ ist eine Konkunktion \\
$L_1 \AND … \AND L_n$ von Literalen $L_i = \begin{cases} A_i \\ \neg A_i \end{cases} , 1 \leq i \leq n$\\
Ein \underline{Maxterm} über $\{A_1, …, A_n\}$ ist eine Disjunktion\\
$L_1 \OR … \OR L_n$ von Literalen $L_i = \begin{cases} A_i \\ \neg A_i \end{cases} , 1 \leq i \leq n$\\

\beispiel{}\\
Sei $\{A_1\}$ Menge von atomaren Formeln\\
Minterme: $A_1$, $\neg A_1$\\
Maxterme: $A_1$, $\neg A_1$\\

\noindent
Sei $\{ A_1, A_2 \}$ Menge von atomaren Formeln\\
Minterme: $(A_1 \AND A_2)$, $(A_1 \AND \neg A_2)$, $(\neg A_1 \AND A_2)$, $\neg ( A_1 \AND \neg A_2)$\\
Maxterme: analog $( \OR )$, $( \OR )$, $( \OR )$, $( \OR )$\\

\lemma{2.1}
Über $\{A_1, …, A_2\}$ gibt es \\
genau $2^n$ Minterme und\\
genau $2^n$ Maxterme.

\beweis{}
Für jedes Literal $L_i$ gibt es (siehe Def.) 2 Möglichkeiten.\\
Also insgesamt $\underbrace{2 \cdot … \cdot 2}_{n-mal} = 2^n$ Möglichkeiten\\

\bemerkung{}
Da $2^n < 2^{2^n}$ für $n \geq 1$ gibt es nicht in jeder semantischen Klasse einen Minterm bzw. Maxterm


\newpage
\datum{28.10.2013}


\noindent
Literale $A_i, \neg A_i$\\
Minterme $L_1 \AND … \AND L_n, L_i = A_i$ oder $\neg A_i$\\
Maxterme $L_1 \OR … \OR L_n, L_i = A_i$ oder $ \neg A_i$\\

\noindent
$2^n$ Minterme, $2^n$ Maxterme über $A_1, …, A_n$\\
$2^n < 2^{2^n} = \# semantische\ Klassen$\\
\noindent\\
Trotzdem wollen wir die semantischen Klassen charakterisieren, in denen MInterme Maxterme liegen.\\

\lemma{2.2}\\
Sei $L_1, …, L_n$ Literale mit $L_i = A_i$ oder $L_i = \neg A_i$\\
für $1 \leq i \leq n$ und $f: \{A_1, …, A_n\} \rightarrow \{W, F\}$ eine Interpretation.\\
Dann gilt\\
$f(L_1 \AND … \AND L_n) = W \leftrightarrow f(A_i) = \begin{cases} W\ f\ddot{u}r\ L_i = A_i \\ F\ f\ddot{u}r\ L_i = \neg A_i \end{cases}$\\
$f(L_1 \OR … \OR L_n) = F \leftrightarrow f(A_i) = \begin{cases} F\ f\ddot{u}r\ L_i = A_i \\ W\ f\ddot{u}r\ L_i = \neg A_i \end{cases}$\\

\beweis{} Nachrechnung

\lemma{2.3}
\begin{enumerate}
\item[(i)] Zu einer Wahrheitstabelle gibt es einen Minterm mit dieser Wahrheitstabelle \\
$\Leftrightarrow$ Nur für eine einzige Interpretation steht W in der letzten Spalte 

\item[(ii)] Zu einer Wahrheitstabelle gibt es einen Maxterm mit dieser Wahrheitstabelle \\
$\Leftrightarrow$ Nur für eine einzige Interpretation steht F in der letzten Spalte
\end{enumerate}

\beweis{}\\
„$\Rightarrow$“ Folgt für (i) und (ii) aus Lemma 2.2\\
„$\Leftarrow$“ 
\begin{enumerate}
\item[(i)] Setzten $L_i = \begin{cases} A_i\ f\ddot{u}r\ f(A_i) = W \\ \neg A_i\ f\ddot{u}r\ f(A_i) = F \end{cases}$\\
für die Interpretation $f$, für die W in der letzten Spalte steht.\\
Dann gilt für beliebige Interpretation $g$ nach Lemma 2.2 \\
$g(L_1 \OR … \OR L_n) = W \Leftrightarrow g(A_i) = \begin{cases} W\ f\ddot{u}r\ L_i = A_i \\ F\ f\ddot{u}r\ L_i = \neg A_i \end{cases}$\\
$\Leftrightarrow f = g$\\
$\Rightarrow$ Der Minterm $L_1 \AND … \AND L_n$ hat die gegebene Wahrheitstabelle

\item[(ii)] Setzen $L_1 = \begin{cases} A_i\ f\ddot{u}r\ f(A_i) = F \\ \neg A_i\ f\ddot{u}r\ f(A_i) = W \end{cases}$\\
für die Interpretation $f$, für die F in der letzten Spalte steht.\\
Dann gilt für beliebige Interpretationen $g$ nach Lemma 2.2\\
$g(L_1 \OR … \OR L_n) = F \Leftrightarrow g(A_i) = \begin{cases} F\ f\ddot{u}r\ L_i = A_i \\ W\ f\ddot{u}r\ L_i = \neg A_i \end{cases}$\\
$\Leftrightarrow g = f$\\
$\Rightarrow$ Der Maxterm $L_1 \OR … \OR$ hat die gegebene Wahrheitstabelle $\Box$

\end{enumerate}

\beispiel{}\\
\begin{tabular}{c|c|c}
$A_1$ & $A_2$ &  \\
\hline
F & F & F \\
F & W & W \\
W & F & W \\
W & W & W \\
\end{tabular}\\
Maxterm $A_1 \OR A_2 $\\
Minterm existiert nicht

\definition{}\\
Eine Formel $G$ über $\{A_1, …, A_n\}$ heißt
\begin{enumerate}
\item[(i)] in disjunktiver Normalform (DNF)\\
falls $G$ eine Disjunktion von Mintermen ist\\
$G = \bigvee_{i=1}^k (L_{i1} \AND … \AND L_{in})$\\
$L_{ij} = \begin{cases} A_j \\ \neg A_j \end{cases} 1 \leq j \leq n$\\

\item[(ii)] in konjunktiver Normalform (KNF)\\
falls $G$ eine Konjunktion von Maxtermen ist\\
$G = \bigwedge_{i=1}^k (L_{in} \OR … \OR L_{in})$\\
$L_{ij} = \begin{cases} A_j \\ \neg A_j \end{cases}$ $1 \leq j \leq n$
\end{enumerate}

\bemerkung{}
\begin{itemize}
\item[$\rightarrow$] Dieser Begriff von KNF, DNF stimmt nicht genau mit dem Begriff aus dem Buch von Schöning überein. Dessen Definition wid bei uns als VDNF, VKNF auftauchen.
\item[$\rightarrow$] Für DNF $\bigvee_{i=1}^k ( \bigwedge_{j=1}^n L_{ij} )$ und \\
KNF $\bigwedge_{i=1}^k ( \bigvee_{j=1}^n L_{ij})$\\
erlauben wir $k=0 \leftarrow$ leere DNF, leere KNF
\end{itemize}
\noindent
Dies sind „keine“ Formeln im strengen Sinn\\
Werden sie aber für Satz 2.4 brauchen\\
Wir setzen\\
die DNF für $k=0$ als unerfüllbar\\
die KNF für $k=0$ als gültig\\

\satz{2.4}
Sei $G$ eine Formel über $\{ A_1, …, A_n\}$\\
Dann gibt es Formeln $G_D$ und $G_K$ mit
\begin{itemize}
\item $G \equiv G_D \equiv G_K$
\item $G_D$ is in DNF\\
$G_K$ ist in KNF
\end{itemize}


\beweis{}
Konstruktion von $G_D$ 
\begin{enumerate}
\item[(i)] Seien $f_1, …, f_k$ die Interpretationen mit $f_i(G) = W$ \hspace{.5cm} $1 \leq i \leq k$
\item[(ii)] Für $1 \leq i \leq k$ konstruiere Minterm $L_{i1} \AND … \AND L_{in}$\\
so dass $f(L_{i1} \AND … \AND L_{in}) = W \Leftrightarrow f = f_i$\\
(Benutze Lemma 2.3)
\item[(iii)] Setze\\
$G_D = \bigvee_{i=1}^k (L_{i1} \AND … \AND L_{in})$ in DNF nach Konstruktion\\
Sei $f$ Interpretation\\
$f(G_D) = W \Leftrightarrow$ Es gibt $1 \leq i \leq k$ mit $f(L_{i1} \AND … \AND L_{in}) = W$\\
$\leftrightarrow$ Es gibt $1 \leq i \leq k$ gilt mit $f(A_j) = \begin{cases} W\ f\ddot{u}r\ L_{ij}=A_j \\ F\ f\ddot{u}r\ L_{ij}=\neg A_j \end{cases}$\\
$\Leftrightarrow$ Es gibt $1 \leq i \leq k$ mit $f = f_i$\\
$\Leftrightarrow f(G) = W$\\
$\Rightarrow G \equiv G_D$\\
\end{enumerate}
\noindent

Konstruktion von $G_K$:
\begin{enumerate}
\item[(i)] Seien $f_1, …, f_k$ die Interpretationen mit $f_i(G) = F$ \hspace{.5cm} $1 \leq i \leq k$
\item[(ii)] Für $1 \leq i \leq k$ konstruiere Maxterm $L_{i1} \OR … \OR L_{in}$\\
so dass $f(L_{in} \OR … \OR L_{in}) = F \Leftrightarrow f = f_i$\\
(Benutze Lemma 2.3)
\item[(iii)] Setze\\
$G_K = \bigwedge_{i=1}^k ( L_{i1} \OR … \OR L_{in})$ in KNF nach Konstruktion\\
Sei $f$ eine Interpretation:\\
$f(G_K) = F \Leftrightarrow$ Es gibt $1 \leq i \leq k$ mit $f(L_{i1} \OR … \OR L_{in}) = F$\\
$\Leftrightarrow^{2.3}$ Es gibt $1 \leq i \leq k$ mit $f(A_j) = \begin{cases} F\ f\ddot{u}r\ L_{ij}= A_j \\ W\ f\ddot{u}r\ L_{ij}= \neg A_j \end{cases}$\\
$\Leftrightarrow$ Es gibt $1 \leq i \leq k$ mit $f = f_i$\\
$\Leftrightarrow f(G) = F$\\
$\Rightarrow G \equiv G_K$ $\Box$
\end{enumerate}

\beispiel{}
$G = A_1 \OR A_2$\\
\begin{tabular}{lc|c|c}
  & $A_1$ & $A_2$ &  \\
\hline
 & F & F & F \\
$f_1=$ & F & W & W \\
$f_2=$ & W & F & W \\
$f_3=$ & W & W & W \\
\end{tabular}\\
\noindent\\
$f_1 \leftrightarrow \neg A_1 \AND A_2$\\
$f_2 \leftrightarrow A_1 \AND \neg A_2$\\
$f_3 \leftrightarrow A_1 \AND A_2$\\
$G_D = ( \neg A_1 \AND A_2 ) \OR (A_1 \neg A_2) \OR ( A_1 \AND A_2)$\\
\noindent\\
Für KNF\\
\begin{tabular}{lc|c|c}
  & $A_1$ & $A_2$ &  \\
\hline
$f_1=$ & F & F & F \\
 & W & F & W \\
 & F & W & W \\
 & W & W & W \\
\end{tabular}

\folgerung{2.5}
Bis auf die Reihenfolge der Minterme und Maxterme und das Setzen der Klammern in einer vollständigen Klammerung gibt es in jeder semantischen Klasse eindeutig bestimmte Formeln in DNF und in KNF\\

\noindent\\
Erfüllbarkeitsproblem:\\
EIngabe Formel $G$:\\
Ausgabe: $G$ erfüllbar oder $G$ unerfüllbar \\
$\rightarrow$ Bruteforce: Feste alle $2^n$ Interpretationen in der $n$ atomaren Formeln aus $G$ \\
$\leadsto$ In worst case $2^n$ Schritte = Laufzeit.\\

\noindent\\
Sei $G$ in DNF\\
Ist $G$ nicht die leere DNF, so ist $G$ erfüllbar.\\
$\leadsto$ Erfüllbarkeit in konstanter Zeit entscheidbar\\

\noindent
Sei $G$ in KNF\\
Hat $G \leq 2^n -1$ Maxterme, so ist $G$ erfüllbar.\\
$\leadsto$ Erfüllbarkeit in linearer Zeit entscheidbar.


\end{document}